\subsubsection{測度とメトリクス}

\term{測度}{Measure} と\term{メトリクス}{Metric} は、設計の大きさ、構造、複雑さ、品質を定量的に見るために使われる。測定対象は、設計方法によって異なる。

機能分解を中心とする設計では、構造チャートなどから、階層の深さ、モジュール数、呼び出し関係の複雑さを測る。オブジェクト指向設計では、クラス図などから、クラス数、継承の深さ、結合度、凝集度、メソッド数、メッセージ数などを測る。

測度は、設計の良し悪しを自動的に決めるものではない。測度は、注意すべき箇所を見つける手がかりである。たとえば、結合度が高い部品は変更の影響が大きい可能性がある。凝集度が低い部品は責務が散らばっている可能性がある。
